\section{Technology Stack}
\label{sec:technology-stack}
% Requires: \usepackage{longtable}, \usepackage{array}

The stack is selected to fit Breakly's small team, limited schedule, device
demonstration and functional requirements. Each layer is compared before the
final choice is presented.

% ---------------------------------------------------------------------
\subsubsection{Mobile Client Framework}

\begingroup
\footnotesize
\begin{longtable}{|>{\raggedright\arraybackslash}p{2.6cm}|>{\raggedright\arraybackslash}p{2.95cm}|>{\raggedright\arraybackslash}p{2.95cm}|>{\raggedright\arraybackslash}p{2.95cm}|>{\raggedright\arraybackslash}p{2.95cm}|}
\caption{Comparison of mobile client frameworks}\label{tab:tech-mobile}\\
\hline
\textbf{Option} & \textbf{Language and ecosystem} & \textbf{Development speed for two people} & \textbf{Testing and demo} & \textbf{Performance and native access} \\ \hline
\textbf{React Native with Expo (selected)} & TypeScript; very large npm ecosystem &

One codebase for Android and iOS; fast refresh; many ready-made modules &

Runs on real phones through Expo Go; APK builds with EAS Build &

Good for list and media screens; native features through Expo modules \\ \hline

Flutter & Dart; smaller but growing ecosystem &

One codebase; fast once Dart is learned &

Needs an emulator or a full build; no Expo Go equivalent &

Excellent rendering; native features through plugins \\ \hline

Native (Kotlin and Swift) & Two languages and two toolchains &

Two codebases; slowest for a team of two &

Separate builds; iOS builds need macOS &

Best performance and full platform access \\ \hline

\end{longtable}

\endgroup

\noindent\textbf{Decision: React Native with Expo (TypeScript) and Expo Router.} A single codebase speeds up development, while Expo Go supports testing on real phones. EAS Build can produce an installable APK. Some native features, such as background audio, may require a development build.

% ---------------------------------------------------------------------

\subsubsection{In-App Architecture and State Management}

\begingroup

\footnotesize

\begin{longtable}{|>{\raggedright\arraybackslash}p{2.6cm}|>{\raggedright\arraybackslash}p{2.95cm}|>{\raggedright\arraybackslash}p{2.95cm}|>{\raggedright\arraybackslash}p{2.95cm}|>{\raggedright\arraybackslash}p{2.95cm}|}

\caption{Comparison of state management approaches}\label{tab:tech-state}\\

\hline

\textbf{Option} & \textbf{Boilerplate and learning curve} & \textbf{Fit for session, theme and player state} & \textbf{Server data handling} & \textbf{Testability} \\ \hline

\textbf{Zustand with TanStack Query (selected)} & Very little boilerplate; small API &

Simple stores that also work outside components, for example the timer &

Query library handles caching, retries and refresh &

Stores are plain functions; easy to unit test \\ \hline

Redux Toolkit with RTK Query & More concepts and files &

Strict structure and good DevTools &

RTK Query covers caching and refresh &

Very testable; heavier than this scope needs \\ \hline

React Context with useReducer & No extra dependency &

Frequent updates such as timer ticks re-render many components &

Caching and retries written by hand &

Testable but manual \\ \hline

\end{longtable}

\endgroup

\noindent\textbf{Decision: Zustand, TanStack Query and layered MVVM.} Views handle presentation, ViewModels (custom hooks and Zustand stores) manage UI state and logic, while repositories encapsulate API, storage and media access. This keeps the app testable without the overhead of full Clean Architecture.

% ---------------------------------------------------------------------

\subsubsection{Backend Framework}

\begingroup

\footnotesize

\begin{longtable}{|>{\raggedright\arraybackslash}p{2.6cm}|>{\raggedright\arraybackslash}p{2.95cm}|>{\raggedright\arraybackslash}p{2.95cm}|>{\raggedright\arraybackslash}p{2.95cm}|>{\raggedright\arraybackslash}p{2.95cm}|}

\caption{Comparison of backend options}\label{tab:tech-backend}\\

\hline

\textbf{Option} & \textbf{Language and team fit} & \textbf{Development speed} & \textbf{Control over API and logic} & \textbf{Runtime and operations} \\ \hline

NestJS (Node.js, TypeScript) & Same language as the app; types can be shared &

Fast; built-in modules, validation and Swagger support &
Full control of RSS aggregation, streak logic and scheduled jobs &
Light on memory; quick to start; easy in Docker \\ \hline

\textbf{Spring Boot (Java) (selected)} & Java; mature ecosystem and strong typing &

Moderate setup; reusable Spring modules and conventions &

Full control; Spring Security and Spring Data JPA &

Heavier on memory and start-up time; runs well in Docker \\ \hline
Backend-as-a-service (Supabase or Firebase) & No server code; vendor-specific rules and SDKs &
Fastest for basic data and authentication &
Limited for custom logic; API mostly generated &
No servers to run; vendor lock-in and quotas \\ \hline
\end{longtable}
\endgroup

\noindent\textbf{Decision: Spring Boot (Java).} Breakly needs custom services for content aggregation, caching, media-source integration and streak calculation. Spring Boot provides a mature modular framework, dependency injection and strong integration with PostgreSQL and security libraries.

% ---------------------------------------------------------------------

\subsubsection{Database}

\begingroup

\footnotesize

\begin{longtable}{|>{\raggedright\arraybackslash}p{2.6cm}|>{\raggedright\arraybackslash}p{2.95cm}|>{\raggedright\arraybackslash}p{2.95cm}|>{\raggedright\arraybackslash}p{2.95cm}|>{\raggedright\arraybackslash}p{2.95cm}|}

\caption{Comparison of database options}\label{tab:tech-db}\\

\hline

\textbf{Option} & \textbf{Fit with Breakly data} & \textbf{Query capabilities} & \textbf{Tooling and ORM support} & \textbf{Hosting and cost} \\ \hline

\textbf{PostgreSQL (selected)} & Relational: users, sessions and preferences linked by foreign keys &

Date aggregation and window functions for streaks; JSONB for flexible fields &

Excellent (Spring Data JPA, Hibernate, migrations) &

Open source; runs in Docker; free \\ \hline

MySQL & Relational; fits the same data &

Good; JSON support is less rich &

Good &

Open source; free \\ \hline

MongoDB & Document store; relations handled in code &

Aggregation pipeline; date logic is less natural &

Good (Mongoose) &

Free community edition; managed tiers have limits \\ \hline

\end{longtable}

\endgroup

\noindent\textbf{Decision: PostgreSQL with Spring Data JPA (Hibernate).} Its relational model fits users, sessions and preferences, while SQL supports date-based streak queries. JSONB accommodates flexible session data, and PostgreSQL runs locally or in Docker.

% ---------------------------------------------------------------------

\subsubsection{Authentication}

\begingroup

\footnotesize

\begin{longtable}{|>{\raggedright\arraybackslash}p{2.6cm}|>{\raggedright\arraybackslash}p{2.95cm}|>{\raggedright\arraybackslash}p{2.95cm}|>{\raggedright\arraybackslash}p{2.95cm}|>{\raggedright\arraybackslash}p{2.95cm}|}

\caption{Comparison of authentication options}\label{tab:tech-auth}\\

\hline

\textbf{Option} & \textbf{Security effort} & \textbf{Control and data ownership} & \textbf{Setup effort} & \textbf{Fit for mobile} \\ \hline

\textbf{Backend-issued JWT with refresh tokens (selected)} & The team must hash passwords and handle expiry and rotation correctly, using well-known libraries &

Users stay in our own PostgreSQL database; no vendor dependency &

Moderate: a few endpoints &

Short-lived access token plus refresh token kept in secure storage \\ \hline

Managed authentication (Firebase Auth, Supabase Auth, Auth0) & Handled by the provider &

Users stored with the provider; vendor lock-in &

Low &

SDKs available; extra dependency \\ \hline

Server-side session cookies & Simple and well understood &

Server keeps session state &

Low &

Cookies are awkward in React Native; tokens are preferred \\ \hline

\end{longtable}

\endgroup

\noindent\textbf{Decision: backend-issued JWT access and rotating refresh tokens.} Passwords are hashed with Argon2id, access tokens are short-lived, and refresh tokens are rotated. The app stores tokens in secure device storage (NFR-04), while user data remains in PostgreSQL.

% ---------------------------------------------------------------------

\subsubsection{API Style}

\begingroup

\footnotesize

\begin{longtable}{|>{\raggedright\arraybackslash}p{2.6cm}|>{\raggedright\arraybackslash}p{2.95cm}|>{\raggedright\arraybackslash}p{2.95cm}|>{\raggedright\arraybackslash}p{2.95cm}|>{\raggedright\arraybackslash}p{2.95cm}|}

\caption{Comparison of API styles}\label{tab:tech-api}\\

\hline

\textbf{Option} & \textbf{Simplicity} & \textbf{Contract and documentation} & \textbf{Mobile fit} & \textbf{Fit for a small API} \\ \hline

\textbf{REST with JSON and OpenAPI (selected)} & Simple and familiar &

OpenAPI gives endpoints, requests, responses and error codes &

Standard HTTP caching; easy to test with Postman &

Ideal for about a dozen endpoints \\ \hline

GraphQL & Needs a schema and resolvers &

Self-describing schema; no standard error-code table &

Precise queries; caching is more complex &

Overkill for a few fixed screens \\ \hline

tRPC & Very low friction in TypeScript &

Shared types, but no standard OpenAPI contract &

Works with TanStack Query &

Weak as a documented contract for evaluators \\ \hline

\end{longtable}

\endgroup

\noindent\textbf{Decision: REST over HTTPS with JSON and OpenAPI.} REST suits the app\textquotesingle s small set of fixed endpoints, while OpenAPI documents requests, responses and error codes as the API contract.

% ---------------------------------------------------------------------

\subsubsection{External Content Sources}

\begingroup

\footnotesize

\begin{longtable}{|>{\raggedright\arraybackslash}p{3.2cm}|>{\raggedright\arraybackslash}p{2.9cm}|>{\raggedright\arraybackslash}p{2.9cm}|>{\raggedright\arraybackslash}p{2.7cm}|>{\raggedright\arraybackslash}p{2.7cm}|}

\caption{Comparison of news sources}\label{tab:tech-news}\\

\hline

\textbf{Option} & \textbf{Terms of use} & \textbf{Content provided} & \textbf{Cost and limits} & \textbf{Reliability} \\ \hline

\textbf{Publisher RSS feeds, fetched and cached by the backend (selected)} & Published for syndication; terms checked per publisher \textbf{[TODO]} &

Title, summary, image and link &

Free; no key &

Depends on each publisher; caching reduces load \\ \hline

News aggregation API (for example NewsAPI or GNews) & Provider terms; free tiers are often limited to development &

Titles, snippets and links &

Free tier with request caps &

Stable, but quotas apply \\ \hline

Scraping article pages & Often against site terms &

Full text possible &

Free but fragile &

Breaks when layouts change \\ \hline

\end{longtable}

\begin{longtable}{|>{\raggedright\arraybackslash}p{3.2cm}|>{\raggedright\arraybackslash}p{2.9cm}|>{\raggedright\arraybackslash}p{2.9cm}|>{\raggedright\arraybackslash}p{2.7cm}|>{\raggedright\arraybackslash}p{2.7cm}|}

\caption{Comparison of music sources}\label{tab:tech-music}\\

\hline

\textbf{Option} & \textbf{Licence} & \textbf{Catalogue and categories} & \textbf{Streaming effort} & \textbf{Reliability} \\ \hline

\textbf{Jamendo API, with a small stored fallback set (selected)} & Creative Commons music; API terms to be verified \textbf{[TODO]} &

Large catalogue with tags such as lo-fi, piano and acoustic &

Streaming URLs and artwork provided &

Depends on an external API; fallback tracks keep demos safe \\ \hline

Audius API & Artist-uploaded; terms vary \textbf{[TODO]} &

Large; relaxing categories less certain &

Streaming URLs provided &

Decentralised network; variable availability \\ \hline

Self-hosted royalty-free tracks & Each licence checked by the team &

Small and fully curated &

The team stores and serves the files &

Fully under our control; most manual work \\ \hline

\end{longtable}

\begin{longtable}{|>{\raggedright\arraybackslash}p{3.2cm}|>{\raggedright\arraybackslash}p{2.9cm}|>{\raggedright\arraybackslash}p{2.9cm}|>{\raggedright\arraybackslash}p{2.7cm}|>{\raggedright\arraybackslash}p{2.7cm}|}

\caption{Comparison of video sources}\label{tab:tech-video}\\

\hline

\textbf{Option} & \textbf{Licence} & \textbf{Content fit} & \textbf{Playback in the app} & \textbf{Reliability and limits} \\ \hline

\textbf{Pexels Videos API (selected)} & Free licence; attribution required \textbf{[TODO: verify]} &

Stock clips (nature, city, travel) that suit a calm theme &

Direct MP4 files play with expo-video &

API key and rate limit; backend caching \\ \hline

YouTube Shorts embed & Embedding allowed under the player terms; controls cannot be altered &

Rich and entertaining, but recommendations can lead to endless content &

WebView embed; limited control; background playback restricted &

Depends on YouTube \\ \hline

Self-hosted clips & Own or licensed clips &

Fully curated &

The team hosts and streams the files &

Under our control; storage and bandwidth cost \\ \hline

\end{longtable}

\endgroup

\noindent\textbf{Decision: publisher RSS, Jamendo API and Pexels Videos API.} The backend fetches and caches all external content, protects API keys and keeps feeds finite. RSS supplies article summaries and links; Jamendo provides music and fallback tracks; Pexels supplies short stock clips. All sources are listed on the credits screen.

% ---------------------------------------------------------------------

\subsubsection{Hosting and Deployment}

\begingroup

\footnotesize

\begin{longtable}{|>{\raggedright\arraybackslash}p{2.6cm}|>{\raggedright\arraybackslash}p{2.95cm}|>{\raggedright\arraybackslash}p{2.95cm}|>{\raggedright\arraybackslash}p{2.95cm}|>{\raggedright\arraybackslash}p{2.95cm}|}

\caption{Comparison of hosting options}\label{tab:tech-hosting}\\

\hline

\textbf{Option} & \textbf{Cost} & \textbf{Setup and operation effort} & \textbf{Control and learning value} & \textbf{Demo reliability} \\ \hline

\textbf{Single VPS with Docker Compose (selected)} & Low, fixed monthly price &

Moderate: one Compose file for the API, PostgreSQL and a reverse proxy &

Full control; easy to reproduce locally &

Always on; one machine is a single point of failure \\ \hline

AWS (EC2 or ECS with RDS) & Higher; free tier is limited &

High: networking, access roles and managed services &

Industry-standard skills &

Highly available when configured; more to configure \\ \hline

Platform-as-a-service (Render, Railway, Fly.io) & Low to medium; free tiers may sleep &

Low: deploy from Git &

Less control &

Cold starts on free tiers can hurt demos \\ \hline

\end{longtable}

\endgroup

\noindent\textbf{Decision: a single VPS with Docker Compose.} One Compose file

starts the API, PostgreSQL and a Caddy reverse proxy that provides HTTPS

certificates automatically. GitHub Actions builds and deploys the stack, and a

nightly database dump provides a backup. A VPS is cheap, always on for

demonstrations, and the same Compose stack can be moved to AWS later if

needed.

% ---------------------------------------------------------------------

\subsubsection{Final Stack Summary}

\begingroup

\small

\begin{longtable}{|>{\raggedright\arraybackslash}p{3.2cm}|>{\raggedright\arraybackslash}p{5.2cm}|>{\raggedright\arraybackslash}p{6cm}|}

\caption{Final technology stack}\label{tab:tech-summary}\\

\hline

\textbf{Layer} & \textbf{Technology} & \textbf{Main reason} \\ \hline

\endfirsthead

\multicolumn{3}{r}{\textit{Continued from previous page}}\\[-2pt]\hline

\textbf{Layer} & \textbf{Technology} & \textbf{Main reason} \\ \hline

\endhead

\multicolumn{3}{r}{\textit{Continue on next page}}\\

\endfoot

\endlastfoot

Mobile client & React Native with Expo (TypeScript), Expo Router &

One codebase, testing on real phones with Expo Go \\ \hline

State and data fetching & Zustand, TanStack Query &

Low boilerplate; built-in caching and retries \\ \hline

Local storage & expo-secure-store (tokens), AsyncStorage (preferences, session

end time, upload queue) &

Secure token storage and simple persistence for the offline behaviour in NFR-03 \\ \hline

Media playback & expo-audio, expo-video &

Maintained Expo modules for audio and video \\ \hline

Backend & Spring Boot (Java) &

Mature Java framework for custom services and integrations \\ \hline

API & REST over HTTPS, JSON, OpenAPI &

Simple, and the OpenAPI document serves as the API contract \\ \hline

Database & PostgreSQL with Spring Data JPA (Hibernate) &

Relational data, date queries for streaks, typed access \\ \hline

Authentication & JWT access and refresh tokens, argon2id hashing &

Users stay in our database; no vendor dependency \\ \hline

News & Publisher RSS feeds via the backend &

Headlines and links within syndication terms \\ \hline

Music & Jamendo API with fallback tracks &

Licensed music with ready-made categories \\ \hline

Short videos & Pexels Videos API &

Licensed clips that play directly in the app \\ \hline

Hosting & VPS, Docker Compose, Caddy &

Low cost, full control, reliable demos \\ \hline

CI/CD and builds & GitHub Actions, EAS Build (APK) &

Automated deployment and an installable build \\ \hline

Testing & JUnit 5, Mockito, Spring Boot Test, React Native Testing Library &

Tests for the Java backend and React Native client \\ \hline

\end{longtable}

\endgroup